\newpage
%%%%%%%%%%%%%%%%%%%%%%%%%%%%%%%%%%%%%%%%%%%%%%%%%%%%%%%%%%%%%%%%
%%%%%%%%%%%%%%%%%%%%%%%%%%%%%%%%%%%%%%%%%%%%%%%%%%%%%%%%%%%%%%%%
%%%%%%%%%%%%%%%%%%%%%%%%%% Enunciado %%%%%%%%%%%%%%%%%%%%%%%%%%%

\begin{myblock}
\phantomsection\addcontentsline{toc}{section}{Parte A | Mistral-7B LoRA}
\section*{Parte A | Mistral-7B LoRA}

    Afinar al menos un modelo Transformer en cada parte (generación y clasificación), eligiéndolo
    de Hugging Face. 

\end{myblock}


%%%%%%%%%%%%%%%%%%%%%%%%%%%%%%%%%%%%%%%%%%%%%%%%%%%%%%%%%%%%%%%%
%%%%%%%%%%%%%%%%%%%%%%%%%%%%%%%%%%%%%%%%%%%%%%%%%%%%%%%%%%%%%%%%

\subsection{Mistral-7B LoRA}

Al ser este el modelo con mejores resultados, creí adecuado colocar un apartado más tecnico respecto
a este transformer. 

Mistral mantiene la arquitectura Transformer Decoder estándar, basada en \textit{Self-Attention}
y \textit{Feed-Forward Networks}, pero introduce modificaciones que reducen complejidad y aumentan
la eficiencia sin pérdida de capa expresiva. 

Recordemos que cada bloque de atención se define como:

\[
    \text{Attention}(Q,K,V) = \text{Softmax} \Big(\frac{QK^T}{\sqrt{d_k}+M}\Big) V
\]

Donde $M$ es la máscara causal que impide que el modelo ``mire a futuro''. Sin embargo, mistral 
modifica el mecanismo con dos adiciones:

\begin{itemize}
    \item \textbf{Sliding Window Attention (SWA)}: en lugar de atender a todos los tokens 
    previos, cada token solo atiende a los últimos $W$ tokens. Por defecto esto es $W = 4096$.
    Esto reduce la complejidad de $O(n^2)$ a $O(nW)$, permitiendo manejar contextos extensos
    con menor saturación de memoria. 
    \item \textbf{Grouped-Query Attention (GQA)}: esto reduce el número de cabezas de clave. En lugar
    de una proyección $(Q,K,V)$ por cabeza, agrupa varias cabezas de query y $h_{kv}$ las de 
    key; entonces $h_q < h_{kv}$. Esto reduce la memoria y el tiempo de inferencia sin sacrificar 
    la expresividad
\end{itemize}

Además, cada bloque Transformer incluye una red feed-forward con activación SiLU (Swish):

\[
    \text{FFN}(x) = W_2 \text{SiLU}(W_1x + b_1) + b_2
\]

Con dimension expansion factor de 4× el tamaño del embedding (similar a LLaMA). La función 
de activación SiLU (Sigmoid Linear Unit), conocdia también como Swish, es útil para Mistral poprque
combina lo mejor de ReLU con la Sigmoide, dando resultado a un entrenamineto más estable y un mejor 
rendimiento general del modelo. Esta se define matemáticamente como:

\[
    \text{SiLU}(x) = x \cdot \sigma(x) = x \cdot \frac{1}{1 + e^{-x}}
\]

A diferencia de ReLU, que tiene una esquina afilada con $x=0$, SiLU es una curva suave y continua. 
Esta suavidad le permite al proceso de optimización ser más estable durante el entrenamiento. De igual
forma se evitan cambios bruscos en los gradientes, permitiendo que el modelo converja de manera
más eficiente. Además, para valores negativos pequeños, la función decrece antes de volver a cero.

Asimismo, Mistral-7B LoRA cuenta con un entrenamiento previo sobre 7.3B de parámetros, y 32 
cabezas de atención. Como optimizador utiliza AdamW + cosine decay.  

Ahora bien, hablando de LoRA (Low-Rank Adaptation), este mecanismo permiote ajustar solo una 
pequeña parte de los parámetros del modelo sin modificar los pesos originales. La idea es que
se descomponga una actualización de $\delta W$ en cada matriz de pesos, transformandose en 
dos matrices de rango bajo:

\[
    W' = W + \delta W = W + BA \;\;\;\;\;\;\; A \in \mathbb{R}^{r \times d}, \; B \in \mathbb{R}^{d \times r}
\]

Donde el rango $r$ es pequeño, para nuestro caso específico fue $r = 16$, lo que implica entrenar
menos del 1\% de los parámetros totales. Esto reduce drásticamente el consumo de VRAM y permite ejecutar 
el fine-tuning completo en una sola GPU Titan RTX de 24 GB, que es lo que teníamos disponible
en el servidor de CIMAT Lab.

\clearpage

\subsection{Resultados}

\begin{lstlisting}[
    style=mystyle,
    caption={Output crudo del modelo Mistral-7B LoRA a nivel de palabra.},
    label={lst:output-mistral-word}
]
In the city lights I find myself
I don't know why
I wanna stay inside tonight
I think it's right
But my heart keeps telling me that I should go
Someplace I know where I belong
So, I take a deep breath, step outside
And follow my heart to where I want to hide

In the shadows of my memories
There's things I can't forget
The times when you and I used to be
Together we had what we wanted yet
A sparkle in our eyes would leave us
To chase dreams that only took us far from here
Now, all those years seem like a blur
And time has passed by, no trace left behind
Like footprints in the sand
My past steps slowly wash away
I see them sinking with the tide
Leaving me with a future I might not know

If you could read my mind
What would you think of who I am?
Have I changed since then?
Did your love bring out the best or the worst in me?
Do you remember how I felt?
\end{lstlisting}
\clearpage

\begin{lstlisting}[
    style=mystyle,
    caption={Output crudo del modelo Mistral-7B LoRA a nivel de palabra.},
    label={lst:output-mistral-word}
]
They say I'm broken but I'm breathing,
I say I'll heal but don't know what I mean
You see they say the heart never forgets,
The pain that you feel just gets passed down to me
And now it's like a game of catch,
Every time I move, the feeling gets passed along

Now we all talk about love and compassion,
But how can we show it when there's nothing left?
When every day feels more desperate than the last,
It's easy to believe what they say

They say, "Jena, there's something wrong with you,
Don't you get it yet? You're already gone,"
So I go looking for answers,
Looking for a cure, anything to be strong

"Jena, you're too broken, you won't get well,"
They say, "Your mind is lost and there's no way back"

Am I too far gone? Did I miss my chance?
Do you think maybe there's hope left?
\end{lstlisting}

\clearpage


\begin{lstlisting}[
    style=mystyle,
    caption={Output crudo del modelo Mistral-7B LoRA a nivel de palabra.},
    label={lst:output-mistral-word}
]
Sometimes my shadow sings louder than I do
A symphony of noise so loud it rings and echoes in the dark of night
It shines with light in its heart, and a deep soulful wisdom
The shadow was a creature of its own right
Sometimes we are just shadows of ourselves
Drowning in our own light as the real us watches on in the background
Someone else.
Some days I wake up and forget that I am me
And everything is a haze and then suddenly you're there
You bring back clarity

My best friend once told me that every person needs their best friend to see them for who they truly are
Every day I watch you and you help me to remember myself
When I look at your face I know the depths of your heart
Your eyes can tell stories and I love how you make me feel
I watch as you go about life, painting a masterpiece
One stroke at a time, coloring outside the lines, sometimes straying completely off course
But you never give up
And when I watch you, I feel like anything
\end{lstlisting}

\subsection{Conclusiones}

El modelo Mistral-7B LoRA mostró un gran desempeño y marcada mejora en la coherencia frente a las
arquitecturas tradicionales recurrentes y el resultado de TinyLLaMA. La pérdida promedio final 
fue de 1.59, lo que se traduce en una perplejidad de entrenamiento de 4.93, mientras que la PPL
de validación se estimó entorno a 5.7. Estras cifras reflejas una genración estable, con poco
sobreajuste, pese a realizarse solo en tres épocas, y el tamaño reducido del corpus. 

A nivel cualitativo, las letras generadas demustran una notable consistencia en cuanto a temática
y carga emocional. El texto del promt ``\textit{In the city lights I find myself}'' se muestra como una
canción con buena estructura lírica y estrofras, cerrando con un tono introspectivo ``\textit{Did your love bring out the best or the worst in me?
Do you remember how I felt?}''.  

La segunda canción,``\textit{They say I’m broken but I’m breathing}'' articula un discurso casi como 
de confesión, creo yo, imitando canciones como Cottonwood, Addict With A Pen o Fake You Out de 
Twenty One Pilots.  Finalmente,``\textit{Sometimes my shadow sings louder than I do}'' destaca por su
tono po´etico y met´aforas visuales, lo que sugiere que Mistral aprendi´o patrones 
estil´ısticos y tem´aticos complejos. Me parece que esto se ve impulsado por un buen promt, que
comparta el estilo y visión de la banda.

En comparación con TinyLLaMA-1.1B, Mistral produce versos más largos y mantiene una coherencia 
global más estable, aunque con una ligera pérdida de espontaneidad léxica. Creo que su ventana
de atención amplia y el uso eficiente de LoRA permitieron modelar dependencias de largo alcance
sin degradar la fluidez. 

Algo que pienso es que, aunque entrenamos con aproximadamente el 1\% de los parámetros con LoRa, el resto del
modelo, i.e. los 7 mil millones de parámetros de Mistral base, siguen determinando el comportamineto
de la distribución linguistica general de los outputs. Es decir, LoRA no borra el conocimiento previo
más bien lo modula a través de modificar un poco las matrices. Creo que por eso, sus salidas, aunque 
muy buenas e interesantes, no logran capturar por completo los clichés o estructuras más específicas
de Twenty One Pilots. Es decir, al final sigue siendo un LLM generalista, que ajusta los sesgos
de atención hacia patrones especídifcos. 

En consecuencia, Mistral se adapta y se combierte en un híbrido de estilo global por su entrenamiento
original y el estilo específico de mi corpus. Recordemos que por descomposición de LoRA: $W' = W + BA$.
Donde $B$ y $A$ son matrices de bajo rango (entrenables con menos del 1\% de los params.). Entonces, $W$
cofica todo el conocimiento general del modelo (con la gramática, semántica y razonamiento previo), y 
$BA$ es un filtro que ajusta a los tokens y define los patrones del modelo sobre los que se dará prioridad.

Se intenta reorientar lo aprendido. LoRA ``dobla'' el espacio de atencion, pero el punto de partida 
sigue siendo el embedding universal de Mistral. Esto, creo, explica por qué el modelo mantiene coherencia 
global y gramatical excelente, pero no siempre puede reproducir con fidelidad total el perfil 
artístico de Twenty One Pilots. Supongo que podríamos hacer más ajustes al finetuning para 
acercarnos cada vez más al estilo de Twenty One Pilots, pero eso se queda como trabajo futuro.


%%%%%%%%%%%%%%%%%%%%%%%%%%%%%%%%%%%%%%%%%%%%%%%%%%%%%%%%%%%%%%%%
%%%%%%%%%%%%%%%%%%%%%%%%%%%%%%%%%%%%%%%%%%%%%%%%%%%%%%%%%%%%%%%%

\clearpage